% Source: https://artofproblemsolving.com/community/c3h3313229_a_lot_of_circles
% Stable identity: aops-topic-3313229-q01
The six hexagon vertices are distance 2 from the origin. The two extra centers lie on opposite perpendicular bisectors of hexagon edges, at distance $2\sqrt3$ from the origin: the edge midpoint is $\sqrt3$ away, and the outward equilateral-triangle height is another $\sqrt3$. Thus a circle centered at the origin of radius $2\sqrt3+1$ contains all nine circles. The two extra centers are $4\sqrt3$ apart, so any enclosing circle must contain two opposite points separated by $4\sqrt3+2$, requiring radius at least $2\sqrt3+1$. Hence the least area is $\boxed{(13+4\sqrt3)\pi}$.
\begin{asy}
import geometry;
usepackage("amsmath");
size(270);
pair O=(0,0);draw(circle(O,1));for(int k=0;k<6;++k)draw(circle(2dir(60*k),1));pair U=2sqrt(3)*dir(30),V=-U;draw(circle(U,1),blue);draw(circle(V,1),blue);draw(circle(O,2sqrt(3)+1),dashed);draw(V--U,dashed);label("$O$",O,S);label("$U$",U,N);label("$V$",V,S);
\end{asy}
